\chapter{Tipus ordinals}\label{cap:ordinals}

\begin{objectius}
\apartat{Prerequisits} Del capítol~\ref{cap:bons-ordres}: bons ordres, segments, rigidesa, comparació per segments i tricotomia. Del capítol~\ref{cap:recursio}: la comparabilitat total. Del capítol~\ref{cap:zorn}: el conjunt de Hartogs \(\Hart(A)\) i el fet que és un bon ordre. Del capítol~\ref{cap:families}: unions indexades de parts. Del capítol~\ref{cap:cantor}: el teorema de Cantor, com a model de Burali--Forti.
\apartat{Objectius} En acabar el capítol, el lector hauria de poder:
\begin{itemize}
  \item definir un tipus ordinal com un bon ordre llevat d'isomorfisme, sense representant canònic, i comprovar la independència del representant de cada definició;
  \item demostrar que els tipus estan totalment ordenats per segments, i ben ordenats dins \(\Hart(A)\);
  \item explicar per què \(\alpha\in\beta\) no és una notació adequada en aquest llibre;
  \item definir zero, successor i límit, i demostrar que són exhaustius, excloents i invariants;
  \item construir el suprem d'una família representada dins un bon ordre com el tipus de la reunió dels segments oberts, i saber per què cal l'acotació;
  \item demostrar el teorema de Burali--Forti i llegir-lo com el paral·lel ordinal del teorema de Cantor.
\end{itemize}
\end{objectius}

Els bons ordres mesuren la \emph{forma} d'un procés que es recorre sense
saltar-se cap pas; però dos bons ordres de la mateixa forma ---isomorfs com a
ordres--- no s'han de distingir. El \emph{tipus ordinal} és, precisament, un
bon ordre considerat llevat d'isomorfisme d'ordre: el que queda quan
s'oblida la naturalesa concreta dels elements i es reté només la manera com
estan ordenats. Aquest capítol desenvolupa els tipus ordinals amb la
disciplina estructural del llibre: es defineixen les nocions
\emph{sobre representants} ---bons ordres concrets--- i després es demostra que
no depenen del representant, gràcies a la unicitat de l'isomorfisme d'ordre
(teorema~\ref{teo:unicitat-iso}). L'ordre entre tipus és la comparació per
segments inicials que ja coneixem; el zero, el successor i els punts límit
s'obtenen de construccions elementals sobre bons ordres; i els suprems es
formen \emph{acotadament}, dins un bon ordre donat, sense parlar mai del
\enquote{conjunt de tots els ordinals} ---que, com veurem amb el teorema de
Burali--Forti, no existeix. El paral·lelisme amb el capítol de Cantor és
deliberat: així com no hi ha una mida màxima, tampoc no hi ha una longitud
d'ordre màxima.

\section{Tipus ordinals i el seu ordre}\label{sec:tipus-ordinals}

\begin{definicio}[Tipus ordinal]\index{Tipus ordinal}\label{def:tipus-ordinal}
Un \emph{tipus ordinal} és un conjunt ben ordenat considerat llevat
d'isomorfisme d'ordre (definició~\ref{def:creixent}). Diem que un bon ordre
\((W,<)\) \emph{representa} el tipus \(\alpha\), i escrivim
\(\alpha=\operatorname{tipus}(W)\), quan \(\alpha\) és la seva classe
d'isomorfisme d'ordre. Dos bons ordres representen el mateix tipus si i només
si són isomorfs com a ordres. \emph{No} escollim cap representant canònic:
treballem amb bons ordres qualssevol i comprovem, cada vegada, que el que
definim no depèn del representant triat.
\end{definicio}

\begin{observacio}[Tipus acotats i teoria global]\label{obs:tipus-acotats}
Com a la comparació de conjunts (observació~\ref{obs:comparacio-metateoria}),
\enquote{tipus ordinal} és, en general, una noció metateòrica: parla de bons
ordres llevat d'isomorfisme, no d'elements d'un conjunt fix. Ara bé, per a un
conjunt \(A\) donat, els tipus dels bons ordres sobre parts d'\(A\) sí que
formen un conjunt ---és el conjunt de Hartogs \(\Hart(A)\)
(definició~\ref{def:hartogs})---, i tota la teoria \emph{acotada} d'aquest
capítol es pot llegir dins un \(\Hart(A)\) prou gran. La distinció importa: els
enunciats sobre \enquote{un tipus ordinal} són de Nivell~I quan es poden
situar dins un \(\Hart(A)\); una teoria \emph{global} de tots els tipus
ordinals alhora seria de Nivell~II i requeriria SR
(observació~\ref{obs:sr}). En aquest capítol ens mantenim, tret d'on
s'indiqui, en el terreny acotat.
\end{observacio}

L'ordre entre tipus és la comparació per segments inicials del
capítol~\ref{cap:bons-ordres}, ara llegida sobre tipus.

\begin{definicio}[Ordre dels tipus ordinals]\index{Ordre dels tipus ordinals}\label{def:ordre-ordinals}
Siguin \(\alpha=\operatorname{tipus}(W)\) i \(\beta=\operatorname{tipus}(V)\).
Definim
\[
\alpha<\beta
\quad\Longleftrightarrow\quad
W\ \text{és isomorf a un segment inicial propi de}\ V
\]
(definició~\ref{def:mes-curt}), i \(\alpha\le\beta\) si \(\alpha<\beta\) o
\(\alpha=\beta\). La unicitat de l'isomorfisme d'ordre
(teorema~\ref{teo:unicitat-iso}) garanteix que la definició \emph{no depèn dels
representants}: si \(W\cong W'\) i \(V\cong V'\), aleshores \(W\) és isomorf a
un segment propi de \(V\) si i només si \(W'\) ho és de \(V'\).
\end{definicio}

\begin{proposicio}[Els tipus ordinals estan totalment ordenats]\label{prop:ordinals-total}
La relació \(<\) sobre tipus ordinals és un ordre estricte total: per a
tipus \(\alpha,\beta\) qualssevol es dona exactament una de \(\alpha<\beta\),
\(\alpha=\beta\), \(\beta<\alpha\). A més, restringida als tipus dels bons
ordres sobre parts d'un conjunt fix \(A\) ---és a dir, sobre \(\Hart(A)\)---,
és un bon ordre.\par
\emph{Axiomes usats:} comparabilitat i tricotomia de bons ordres
(teoremes~\ref{teo:comparabilitat} i~\ref{teo:tricotomia}), Hartogs
(proposició~\ref{prop:hartogs-bon-ordre}); cap axioma nou; \emph{cap elecció}.
\end{proposicio}

\begin{prova}
La tricotomia i la comparabilitat dels bons ordres
(teoremes~\ref{teo:tricotomia} i~\ref{teo:comparabilitat}) diuen exactament
que, per a dos representants \(W\) i \(V\), es dona una i una sola de: \(W\)
isomorf a un segment propi de \(V\) (és a dir \(\alpha<\beta\)), \(W\cong V\)
(és a dir \(\alpha=\beta\)), o \(V\) isomorf a un segment propi de \(W\) (és a
dir \(\beta<\alpha\)). La transitivitat és la de la comparació per segments
(proposició~\ref{prop:comparacio-ordres}), i la irreflexivitat és que cap bon
ordre no és isomorf a un segment propi seu
(corol·lari~\ref{cor:no-segment-propi}). Que sobre \(\Hart(A)\) és un bon
ordre és la proposició~\ref{prop:hartogs-bon-ordre}.
\end{prova}

\begin{observacio}[No s'ha d'escriure \(\alpha\in\beta\)]\label{obs:no-pertinença-ordinal}
Hi ha una diferència important amb la teoria clàssica de von Neumann, on un
ordinal \emph{és} el conjunt dels ordinals menors i \enquote{\(\alpha<\beta\)}
es defineix com \enquote{\(\alpha\in\beta\)}. En aquest llibre la pertinença
és sempre \emph{relativa a una part d'un conjunt fix}
(capítol~\ref{cap:parts-interna}), i no hi ha cap relació de pertinença global
entre conjunts; per tant \(\alpha\in\beta\) \emph{no} és una notació adequada
per a l'ordre ordinal. La relació correcta és \(\alpha<\beta\), definida
mitjançant segments inicials (definició~\ref{def:ordre-ordinals}). Aquesta no
és una qüestió de gust: escriure \(\alpha\in\beta\) suggeriria una jerarquia
material de pertinença que el llibre no assumeix, i que precisament el punt de
vista estructural vol evitar.
\end{observacio}

\section{Zero, successor i límit}\label{sec:zero-successor-limit}

Les tres classes bàsiques de tipus ordinals es defineixen sobre
representants, amb construccions elementals de bons ordres, i després es
comprova la invariància.

\begin{definicio}[Zero i successor]\index{Zero i successor}\label{def:zero-successor}
El \emph{tipus zero}, \(0\), és el tipus del bon ordre \emph{inicial}
\(\mathbf0\) (l'ordre buit): el menor de tots els tipus, perquè \(\mathbf0\) és
isomorf a un segment inicial de qualsevol bon ordre. Donat un bon ordre
\((W,<)\), el seu \emph{successor} és el tipus del bon ordre \(W+\mathbf1\)
---la suma disjunta (definició~\ref{def:suma-disjunta}) amb un punt nou
\(\ast\) col·locat \emph{després} de tots els elements de \(W\): l'ordre en
què \(u<\ast\) per a tot \(u\in W\) i la restricció a \(W\) és l'original. Si
\(\alpha=\operatorname{tipus}(W)\), el successor es denota \(\alpha+1\).
\end{definicio}

\begin{proposicio}[Invariància i caracterització del successor]\label{prop:successor-invariant}
El successor no depèn del representant: si \(W\cong V\), aleshores
\(W+\mathbf1\cong V+\mathbf1\). A més, per a tot bon ordre \(W\):
\begin{enumerate}
\item \(\alpha<\alpha+1\), i \(\alpha+1\) és el \emph{menor} tipus
estrictament més gran que \(\alpha\): no hi ha cap tipus \(\gamma\) amb
\(\alpha<\gamma<\alpha+1\);
\item \(\alpha+1\) té un element \emph{màxim} (el punt nou \(\ast\)), i el
segment inicial que aquest determina és, precisament, \(\alpha\).
\end{enumerate}
\emph{Axiomes usats:} suma disjunta (capítol~\ref{cap:valors}), unicitat de
l'isomorfisme d'ordre (teorema~\ref{teo:unicitat-iso}); cap elecció.
\end{proposicio}

\begin{prova}
\emph{Invariància.} Un isomorfisme d'ordre \(f\colon W\cong V\) s'estén a
\(f+\id\colon W+\mathbf1\cong V+\mathbf1\) enviant \(\ast\) a \(\ast\); conserva
l'ordre perquè els punts nous ocupen la mateixa posició final. \emph{(1)} El
segment inicial de \(W+\mathbf1\) determinat pel punt \(\ast\) és exactament
\(W\), de manera que \(W\) és isomorf a un segment propi de \(W+\mathbf1\), és a
dir \(\alpha<\alpha+1\). Si hi hagués \(\gamma\) amb \(\alpha<\gamma<\alpha+1\),
un representant \(G\) de \(\gamma\) seria isomorf a un segment propi de
\(W+\mathbf1\); els segments propis de \(W+\mathbf1\) són \(W\) mateix i els
segments de \(W\) (proposició~\ref{prop:parts-inicials}), de manera que
\(\gamma\le\alpha\), contra \(\alpha<\gamma\). \emph{(2)} El punt \(\ast\) és
màxim per construcció, i el seu segment inicial és \(W\), que representa
\(\alpha\) (proposició~\ref{prop:segment-determina-punt}).
\end{prova}

\begin{definicio}[Tipus límit]\index{Tipus límit}\label{def:tipus-limit}
Un tipus ordinal \(\lambda\) és \emph{successor} si \(\lambda=\alpha+1\) per a
algun \(\alpha\); equivalentment, si un (tot) representant seu té un element
màxim. És \emph{límit} si no és zero ni successor: és a dir, si és no inicial i
cap representant seu no té element màxim ---per a tot element hi ha un de
posterior.
\end{definicio}

\begin{proposicio}[Zero, successor i límit són exhaustius i invariants]\label{prop:tricotomia-successor}
Tot tipus ordinal és exactament un de: zero, successor o límit; i la classe a
què pertany no depèn del representant.\par
\emph{Axiomes usats:} bon ordre (definició~\ref{def:bon-ordre}), unicitat de
l'isomorfisme d'ordre (teorema~\ref{teo:unicitat-iso}); cap elecció.
\end{proposicio}

\begin{prova}
Sigui \(W\) un representant de \(\lambda\). Si \(W\) és inicial, \(\lambda=0\).
Si no, \(W\) té un element mínim però pot tenir o no un element màxim. Tenir
element màxim és una propietat invariant per isomorfisme d'ordre (un
isomorfisme envia màxim a màxim), de manera que no depèn del representant. Si
\(W\) té màxim \(m\), aleshores \(W\cong W_{<m}+\mathbf1\) ---el segment
\(W_{<m}\) seguit de \(m\)---, de manera que \(\lambda=\operatorname{tipus}(W_{<m})+1\)
és successor. Si \(W\) no té màxim i no és inicial, \(\lambda\) no és zero (no
inicial) ni successor (un successor té màxim, proposició~\ref{prop:successor-invariant}),
i per tant és límit. Les tres classes són mútuament excloents ---zero és
inicial, successor té màxim, límit no--- i exhaustives.
\end{prova}

\begin{observacio}[Lectura clàssica; el primer límit]\label{obs:limit-classic}
En la notació habitual, \(0,1,2,\dots\) són els tipus dels bons ordres finits
(els \(\seg{n}\), capítol~\ref{cap:finits}), tots zero o successors; el primer
tipus \emph{límit} és \(\omega=\operatorname{tipus}(\N)\), el tipus del bon
ordre dels naturals (teorema~\ref{teo:bon-ordre-N}), que no té element màxim.
Després venen \(\omega+1,\omega+2,\dots\), successors, i el següent límit
\(\omega+\omega=\omega\cdot2\), i així successivament. La distinció
successor/límit és la que governa les definicions per recursió transfinita:
als successors es fa un pas, als límits es pren un suprem
(capítol~\ref{cap:recursio}). És per als límits, on no hi ha \enquote{element
anterior}, que la recursió ha de mirar tota la història i no només el terme
previ.
\end{observacio}

\section{Suprems acotats}\label{sec:suprems}

Una família de tipus ordinals, si està \emph{representada dins un bon ordre}
comú, té suprem; i el suprem es forma acotadament, sense sortir d'aquell bon
ordre. Aquesta és la manera estructural de reunir tipus sense invocar el
\enquote{conjunt de tots els ordinals}.

\begin{definicio}[Família de tipus representada dins un bon ordre]\index{Família de tipus representada dins un bon ordre}\label{def:familia-tipus}
Una família de tipus ordinals \((\alpha_i)_{i\in I}\) està \emph{representada
dins} un bon ordre \(W\) si existeix una aplicació \(r\colon I\to W\) tal que,
per a cada \(i\), el tipus \(\alpha_i\) és el del segment inicial
\(W_{<r(i)}\). És a dir, cada \(\alpha_i\) apareix com el tipus d'un segment de
\(W\), i la família és la imatge d'una aplicació al conjunt d'índexs
(capítol~\ref{cap:families}).
\end{definicio}

\begin{proposicio}[Suprem acotat]\label{prop:suprem}
Sigui \((\alpha_i)_{i\in I}\) una família de tipus representada dins un bon
ordre \(W\) per \(r\colon I\to W\), és a dir, amb
\(\alpha_i=\operatorname{tipus}(W_{<r(i)})\). Aleshores existeix el
\emph{suprem} \(\sup_i\alpha_i\): el menor tipus ordinal \(\ge\) que tots els
\(\alpha_i\). Concretament, és el tipus de la reunió
\[
S\;=\;\bigcup_{i\in I}W_{<r(i)}
\]
dels segments inicials \emph{oberts}, que és una part inicial de \(W\). Si
\(I\) és inicial, \(S\) és buida i el suprem és \(0\).\par
\emph{Axiomes usats:} unions indexades de parts (capítol~\ref{cap:families}),
bon ordre i segments (capítol~\ref{cap:bons-ordres}); cap axioma nou; \emph{cap
elecció}.
\end{proposicio}

\begin{prova}
Sigui \(S=\bigcup_{i\in I}W_{<r(i)}\hookrightarrow W\) la unió indexada
(capítol~\ref{cap:families}). Com a unió de parts inicials, \(S\) és una part
inicial de \(W\), i per tant un segment \(W_{<w_0}\) o bé tot \(W\)
(proposició~\ref{prop:parts-inicials}); en qualsevol cas és ben ordenat
(proposició~\ref{prop:segment-ben-ordenat}). Sigui
\(\sigma=\operatorname{tipus}(S)\).

\emph{\(\sigma\) és cota superior.} Cada \(W_{<r(i)}\) és una part inicial de
\(S\) (és inicial a \(W\) i està continguda a \(S\)), de manera que és tot
\(S\) o un segment propi de \(S\); en tots dos casos \(\alpha_i\le\sigma\).

\emph{\(\sigma\) és la menor.} Sigui \(\beta\) amb \(\beta\ge\alpha_i\) per a
tot \(i\), amb representant \(V\). Per a cada \(i\) hi ha un únic
isomorfisme d'ordre \(f_i\) de \(W_{<r(i)}\) sobre una part inicial de \(V\)
(teorema~\ref{teo:unicitat-iso}). Aquests isomorfismes són compatibles: si
\(r(i)\le r(j)\), aleshores \(W_{<r(i)}\) és una part inicial de
\(W_{<r(j)}\), la restricció de \(f_j\) és un isomorfisme d'ordre de
\(W_{<r(i)}\) sobre una part inicial de \(V\) (la imatge d'una part inicial
per un isomorfisme sobre una part inicial és inicial), i per la unicitat
coincideix amb \(f_i\); i per la totalitat de \(W\), dos índexs sempre són
en un dels dos casos. La reunió dels grafs de les \(f_i\) és, doncs, una
relació funcional i total sobre \(S\) (proposició~\ref{prop:aplicacio-funcional}),
és a dir, una aplicació \(f\colon S\to V\) que és creixent, injectiva i té
imatge inicial, perquè cada punt de \(S\) i cada desigualtat entre dos punts
de \(S\) viuen dins un mateix \(W_{<r(j)}\), on \(f\) és \(f_j\). Així \(S\)
és isomorf a una part inicial de \(V\), i \(\sigma\le\beta\). Per tant
\(\sigma=\sup_i\alpha_i\).
\end{prova}

\begin{observacio}[Per què els segments oberts]\label{obs:suprem-oberts}
La reunió dels segments \emph{tancats} \(W_{\le r(i)}\) dona en general un
tipus massa gran: si \(I=\mathbf1\) i \(r\) és el mínim de \(W\), aleshores
\(\alpha=\operatorname{tipus}(W_{<r})=0\), mentre que \(W_{\le r}\) té un
element i tipus \(1\). El suprem d'una família amb màxim és aquest màxim, i
la fórmula amb segments oberts ho respecta: si \(r(i_0)\) és el màxim dels
\(r(i)\), aleshores \(S=W_{<r(i_0)}\) i \(\sigma=\alpha_{i_0}\).
\end{observacio}

\begin{observacio}[Per què el suprem ha de ser acotat]\label{obs:suprem-acotat}
La hipòtesi \enquote{representada dins un bon ordre \(W\)} és essencial i és el
que manté la construcció de Nivell~I. Reunir una família \emph{arbitrària} de
tipus ordinals ---sense un \(W\) comú que els contingui tots com a segments---
no es pot fer amb els axiomes actuals: caldria formar primer un bon ordre prou
llarg, i si els tipus no estan acotats cap axioma de Nivell~I no en garanteix
cap; i si la família és la de \emph{tots} els tipus, un tal bon ordre és
impossible, que és el contingut del teorema de Burali--Forti que ve tot
seguit. Per això, cada vegada que prenem un suprem, hem d'exhibir abans el bon
ordre \(W\) dins el qual la família viu ---sovint un \(\Hart(A)\) prou gran. És
la mateixa disciplina de la guia: no escriure \(\sup_i\alpha_i\) sense haver
construït primer la família acotada que es reuneix.
\end{observacio}

\section{El teorema de Burali--Forti}\label{sec:burali-forti}

Així com el teorema de Cantor prohibeix una mida màxima, el teorema de
Burali--Forti prohibeix un \enquote{conjunt de tots els tipus ordinals}: la
col·lecció de tots ells, si fos un conjunt ben ordenat, tindria un tipus
estrictament menor que ell mateix.

\begin{teorema}[Burali--Forti]\label{teo:burali-forti}
No existeix cap conjunt \(\Omega\) que contingui un representant de
\emph{cada} tipus ordinal i que estigui ben ordenat per la relació de segment
inicial de manera que el tipus del seu segment \(\Omega_{<\gamma}\) sigui,
per a cada \(\gamma\in\Omega\), el tipus representat per \(\gamma\).\par
\emph{Axiomes usats:} bon ordre i segments, no-isomorfisme amb segments propis
(corol·lari~\ref{cor:no-segment-propi}); cap axioma nou; \emph{cap elecció}.
\end{teorema}

\begin{prova}
Suposem que un tal \(\Omega\) existís. Aleshores \((\Omega,<)\) seria un bon
ordre, i tindria ell mateix un tipus \(\omega_\Omega=\operatorname{tipus}(\Omega)\).
Com que \(\Omega\) conté un representant de \emph{cada} tipus ordinal, conté un
element \(\gamma\) que representa \(\omega_\Omega\): és a dir, per la hipòtesi
sobre \(\Omega\), \(\operatorname{tipus}(\Omega_{<\gamma})=\omega_\Omega=\operatorname{tipus}(\Omega)\).
Però \(\Omega_{<\gamma}\) és un segment inicial \emph{propi} de \(\Omega\)
(propi perquè \(\gamma\in\Omega\)), i acabem de dir que és isomorf a \(\Omega\)
sencer. Això contradiu que cap bon ordre no és isomorf a cap segment inicial
propi seu (corol·lari~\ref{cor:no-segment-propi}). Per tant no existeix cap tal
\(\Omega\).
\end{prova}

\begin{observacio}[Què diu exactament el teorema]\label{obs:burali-hipotesi}
L'enunciat és precís i convé no ampliar-lo en llegir-lo informalment. Afirma
que no hi ha cap conjunt \(\Omega\) que \emph{alhora} (i)~contingui un
representant de cada tipus, (ii)~estigui ben ordenat per segments, i
(iii)~tingui la propietat que el segment \(\Omega_{<\gamma}\) sigui del tipus
que \(\gamma\) representa. La frase \enquote{no hi ha conjunt de tots els
ordinals} n'és la lectura habitual, i és correcta en aquest sentit: no hi ha
cap conjunt ben ordenat que \emph{indexi coherentment} tots els tipus. El que
el teorema no diu és res sobre col·leccions de representants sense aquesta
estructura d'ordre, que en aquest llibre ni tan sols es poden formular sense
un conjunt ambient. Com que els tipus no són elements de cap conjunt fix
(observació~\ref{obs:tipus-acotats}), la condició (iii) és la manera de dir
\enquote{\(\Omega\) és el conjunt dels tipus} sense pertinença global.
\end{observacio}

\begin{observacio}[Burali--Forti com a paral·lel de Cantor]\label{obs:burali-cantor}
El teorema és l'anàleg ordinal exacte del teorema de Cantor
(teorema~\ref{teo:cantor}), i comparteix amb ell l'estructura diagonal: la
contradicció surt de fer que la construcció s'apliqui a si mateixa. Cantor:
\(A\prec P(A)\), de manera que no hi ha conjunt universal ni mida màxima
(corol·lari~\ref{cor:no-universal}). Burali--Forti: no hi ha un bon ordre de
tots els tipus ordinals, de manera que no hi ha longitud d'ordre màxima. En
tots dos casos, la moralitat és que la col·lecció de \enquote{tots els}
objectes d'una mena ---totes les parts, tots els tipus--- és massa gran per
ser un conjunt del mateix tipus. És per això que la teoria \emph{global} dels
ordinals és de Nivell~II (observació~\ref{obs:tipus-acotats}): reunir tots els
tipus demanaria un principi ---SR--- que els axiomes de Nivell~I no donen, i
que ni tan sols amb SR produiria un \emph{conjunt} ben ordenat de tots ells,
sinó una família estratificada.
\end{observacio}

\begin{resumcap}
\apartat{Idees centrals}
\begin{enumerate}
  \item Un tipus ordinal és la classe d'un bon ordre llevat d'isomorfisme d'ordre; un bon ordre concret de la classe n'és un \emph{representant}, i el seu conjunt subjacent, el \emph{suport}. No hi ha representant canònic, i cada definició es comprova independent del representant gràcies a la unicitat de l'isomorfisme d'ordre.
  \item \(\alpha<\beta\) vol dir que un representant de \(\alpha\) és isomorf a un segment propi d'un de \(\beta\); l'ordre és total, i dins \(\Hart(A)\) és un bon ordre. No s'escriu \(\alpha\in\beta\): no hi ha pertinença global.
  \item El successor \(\alpha+1\) és el tipus de \(W+\mathbf1\), el menor tipus més gran que \(\alpha\). Tot tipus és exactament zero (representant buit), successor (representant amb màxim) o límit (representant no buit i sense màxim); \(\omega=\operatorname{tipus}(\N)\) és el primer límit.
  \item El suprem d'una família representada dins \(W\) per \(r\colon I\to W\) és el tipus de \(\bigcup_iW_{<r(i)}\), la reunió dels segments \emph{oberts}; l'acotació és imprescindible.
  \item Burali--Forti: no hi ha cap bon ordre \(\Omega\) que indexi coherentment tots els tipus (cada \(\gamma\in\Omega\) representa el tipus del segment \(\Omega_{<\gamma}\)). És el paral·lel ordinal de Cantor, i és per això que la teoria global dels ordinals queda fora del Nivell~I.
\end{enumerate}
\apartat{Errors típics} Escollir un representant canònic i oblidar comprovar la independència; escriure \(\alpha\in\beta\); confondre \(\mathbf1+W\) (punt nou al principi, mateix tipus que \(W\) si \(W\) és infinit sense màxim, com \(\mathbf1+\N\cong\N\)) amb \(W+\mathbf1\); prendre el suprem amb segments tancats \(W_{\le r(i)}\), que dona un tipus massa gran; i parlar del \enquote{conjunt de tots els ordinals}.
\apartat{Lectura recomanada} \cite[§§19--21]{halmos} per als ordinals, la seva aritmètica inicial i el teorema de Burali--Forti en la presentació de von Neumann; \cite[cap.~2]{jech2003} i \cite[cap.~I.7--I.8]{kunen2011} per a la teoria clàssica dels ordinals; i \cite[§§8.4--8.5]{lawvere-rosebrugh} per a la lectura estructural dels bons ordres com a tipus.
\end{resumcap}
